\phantomsection\bheading{section}{5.\quad The hard cases}{section}{\protect\numberline{5.}The hard cases}

Each case is stated, run through the rules, and marked \textbf{solved}, \textbf{solved with a cost}, or \textbf{open}.

\phantomsection\bheading{subsection}{5.1\quad Closures that only read a captured list; closures called many times}{subsection}{\protect\numberline{5.1}Closures that only read a captured list; closures called many times}

\begin{Verbatim}
total = List.foldl xs 0 (λx acc → acc + List.length ys)     -- reads ys, called n times
zs = List.set ys 0 7                                        -- then edits ys
\end{Verbatim}

The lambda's summary: $\mathit{cap}\ \mathit{ys}=\mathsf{borrowed}$ (it only reads), $\mathit{once}=\mathsf{no}$. \code{foldl}'s class for its callback: $\mathit{calls}\ \mathsf{many},\allowbreak{}\mathit{escapes}\ \mathsf{no}$. The edge check passes (a non-$\mathit{once}$ closure may be called many times). The closure is a holder of \code{ys} while the closure value is live, which is during the \code{foldl} call only ($\mathit{escapes}=\mathsf{no}$), so at \code{List.set ys 0 7} the holder is dead. \textbf{Solved}, and the reason it is solved is that beni's lambdas are in front of the checker: this is the case research 69 §1.3 item 9 called “the open research question that matters most”, open in every type-based system because a capture counts as a use of the capture's \emph{type} (Wansbrough's $\vdash$-Abs: “all free variables of an abstraction have at least the usage of the abstraction itself”, 1999 §6.1 l. 294; Clean's \code{addInPlace} false error, research 69 §2.1), and closed in the effect- and capture-based ones only with declarations (Deng et al. 2025 Fig. 6; Capybara ms.tex:647–671; Affe's explicit shared borrow \code{\&x}, Radanne et al. 2020 §3.3).

\begin{Verbatim}
ys2 = List.map xs (λx → List.push acc x)                     -- consumes a capture, called n times
\end{Verbatim}

The lambda's summary: $\mathit{cap}\ \mathit{acc}=\mathsf{consumed}$, $\mathit{once}=\mathsf{yes}$. \code{map}'s class: $\mathit{calls}\ \mathsf{many}$. The edge $\mathit{once}\Longrightarrow \mathit{calls}\le 1$ fails: \textbf{rejected}, correctly — the first call writes \code{acc} and returns it, the second writes the same array, so the result list's first element would change under the second call. The message names the lambda, the capture and \code{map}'s call count, and offers the fold (§6.2, example 4). This is OxCaml's lock rule, “a closure that consumes a captured value uniquely must be $\mathit{once}$” (Lorenzen et al. 2024 §3.3, p. 10), and Futhark's refusal of a free- variable consume (2022 post). A closure that consumes a capture and \emph{is} called at most once — \code{Task.\allowbreak{}scope (λscope → … List.\allowbreak{}push acc x …)}, \code{Result.\allowbreak{}andThen r (λx → List.\allowbreak{}push acc x)} — passes, because those callees' classes say $\mathit{calls}\le 1$. \textbf{Solved}, with the cost that a $\mathit{once}$ closure passed to a $\mathsf{many}$ position is an error even when the program would only ever call it once for reasons the checker cannot see (a \code{List.map} over a one-element list).

\phantomsection\bheading{subsection}{5.2\quad \code{▷} pipelines, \code{identity}, composition and polymorphic plumbing}{subsection}{\protect\numberline{5.2}\code{▷} pipelines, \code{identity}, composition and polymorphic plumbing}

\code{▷} is syntax: \code{xs ▷ List.\allowbreak{}filter p ▷ List.\allowbreak{}set 0 v} \emph{is} \code{List.\allowbreak{}set (List.\allowbreak{}filter xs p) 0 v} (\code{language.md} §6.7: “pipes rewrite before anything runs”), and a \code{\_\allowbreak{}} placeholder or a \code{<-} bind is a written lambda (\code{language.md} §6.7), so there is no plumbing function to lose precision in. \code{identity x = x} is inferred $\mathit{res}=\{\pi _{1}\}$; a user-written \code{apply f x = f x} is inferred $\mathit{res}=g.\mathit{res}$ through its class; a composition \code{λx → g (f x)} likewise. This is the case Futhark called “lossy” (2026 aliasing post, ll. 555–612) and proposed to fix by parametricity; here the body is analysed and the class carries the result's provenance, so \code{id}, \code{apply} and \code{|>} keep exact aliasing with no free theorem needed. \textbf{Solved.}

The residual cost is the identity promises of \code{backend.md} §4 (\emph{Identity: what an operation returns unchanged}): \code{filter} that keeps every element returns \code{xs} itself, \code{map} whose results are all \code{===} returns \code{xs}, \code{slice} of everything returns \code{xs}. Under the rule those are $\mathit{res}=\{\mathsf{fresh},\allowbreak{}\pi _{1}\}$ — the result \emph{may} be the input — so \code{List.\allowbreak{}set (List.\allowbreak{}filter xs p) 0 v} demands \code{xs} unique too. In a pipeline \code{xs} is usually dead after, so it costs nothing; where \code{xs} is used later it is a rejection whose message names the promise. §10 (change 2) asks the owner whether the promises stay.

\phantomsection\bheading{subsection}{5.3\quad Generic code under \code{where} clauses}{subsection}{\protect\numberline{5.3}Generic code under \code{where} clauses}

\code{where k.\allowbreak{}compare : k, k → Order} and \code{a.eq} are \code{sync}, read-only by contract: $\mathsf{borrowed}$. A user method's summary is a class on the constraint's method type, bound at each site from the dispatch table's term (§3.7). Inside the generic body, a call \code{x.m a} is a call through a function-typed value whose class is the constraint's; the body's own summary is expressed over it. \textbf{Solved} for precision; the only limit is that a generic body cannot write a value of type $a$ in place (there is no demand of type \code{a → a}), which is the right answer.

\phantomsection\bheading{subsection}{5.4\quad Lists in containers}{subsection}{\protect\numberline{5.4}Lists in containers}

\textbf{A record.} \code{\{ m | rows = List.\allowbreak{}push m.\allowbreak{}rows x \}}: the operand is $(m,\allowbreak{}.\mathit{rows})$; the holders of that cell are $(m,\allowbreak{}.\mathit{rows})$ itself and whatever else holds $m$ or its \code{rows}. The update's other fields are the very values (\code{language.md} §11.12), so the new record and $m$ share every \emph{other} cell — but not \code{rows}, which the push's result replaces. After the update $m$ is dead unless read later; a later \code{m.name} reads a path that does not meet $.\mathit{rows}$, so it is not a use of the cell (§2.7's per-path liveness). \textbf{Solved}: a model threaded through helpers that each touch one field is accepted, because $\mathit{use}$ and \code{live} are per path, which is what the TEA shape needs (§5.9) and what research 42's O8 found necessary.

\textbf{A tuple, a constructor.} The same, through $.i$ and $C\#i$.

\textbf{A list of lists.} \code{List.\allowbreak{}update rows i (λr → \{ r | tags = List.\allowbreak{}push r.\allowbreak{}tags t \})}: the lambda's parameter $r$ is an element of \code{rows}. Two slots of \code{rows} may hold the same record (\code{[ r, r ]}), in which case a push through slot $i$ would be seen at the other slot; the cell-set domain cannot tell slots apart (§2.9). So the derivation needs two things: \code{List.update} must be a consuming writer that overwrites slot $i$ with the callback's result — a destructive read, so the slot's old holder is dead during the callback — and \code{rows} must be \textbf{unique and exclusive}. Then $r$ is unique inside the callback, \code{r.tags}'s only holder is $r$, and the push is in place: accepted. With \code{rows} unique but not exclusive (its elements came from a \code{foreign} that declares nothing, or from \code{List.repeat}), it is rejected as a \textbf{limit} — “the elements of \code{rows} may share” — with the fix \code{List.copy r.tags}. With \code{map} in place of \code{update} under the rule as asked, \code{map} is not a writer and $\mathit{supply}=\mathsf{borrowed}$: the old \code{rows} holds every element until \code{map} has built its result, so the push is \textbf{rejected} as real sharing; under §10's change 1 \code{map} becomes a destructive read per position and the exclusive case is accepted. This is Lean's \code{groupBy}, Clean's propagation rule and the two-slots-one-record hazard in one example. \textbf{Solved with a cost}: a nested edit needs exclusivity, which the program's own construction of the container must establish, and which §9 counts.

\textbf{A \code{Dict}.} \code{Dict k v} is a tree of constructors (\code{static-dispatch-spike.\allowbreak{}md} §5.3), written in beni. \code{Dict.get d k} is a non-destructive read: its result aliases a value inside $d$ ($\mathit{res}=\{\pi _{1}\ \mathit{contents}\}$), and $d$'s values stop being exclusive while that result is live. \code{Dict.\allowbreak{}insert d k v} stores $v$ ($\mathit{use}(v)=\mathsf{shared}$), rebuilds a path and shares the rest with $d$ ($\mathit{res}\ni \pi _{1}$); it keeps $\mathit{excl}$ only if $v$ was unique at the store. \code{Dict.\allowbreak{}update d k f}, written as a path-copying walk that takes the node's value, calls $f$ and rebuilds the node with the result, is a destructive read: its callback's argument is unique iff $d$ is consumed at the call and $d$'s values are exclusive. The Lean counter-example, \code{let group = RBMap.\allowbreak{}find? result x; … RBMap.\allowbreak{}insert result x (Array.\allowbreak{}push group x)} (Huisinga 2023 §2.1, p. 7), is \textbf{rejected} as real sharing with the holder named — \code{result} is live across the push, it is passed to \code{insert} after — and the message offers \code{Dict.update}, which is accepted when the values are exclusive (B.4). \textbf{Solved with the same cost.} Lean accepts the \code{find?} form at run time through its reference count and silently copies; this checker rejects it with the fix on the message.

\phantomsection\bheading{subsection}{5.5\quad Recursion and folds}{subsection}{\protect\numberline{5.5}Recursion and folds}

\begin{Verbatim}
go xs acc = case xs of
    [] → acc
    [ x, …rest ] → go rest (List.push acc (f x))
\end{Verbatim}

\code{go}'s group is one declaration. Optimistic start: $\mathit{acc}\ \mathsf{borrowed},\allowbreak{}\mathit{res}\{\mathsf{fresh}\}$. Round 1: the tail call passes \code{List.push acc …} — a demand on \code{acc}'s cell; holders: \code{acc} (the operand), and nothing else in the frame; after the jump \code{acc} is rebound, so dead. $\mathit{use}(\mathit{acc})=\mathsf{consumed}$; $\mathit{res}=\{\pi _{2}\}$ (the \code{[]} arm returns \code{acc}) $\cup$ the recursive call's $\mathit{res}$, which under the round's assumption is $\{\mathsf{fresh}\}$; round 2 with $\mathit{res}=\{\pi _{2},\allowbreak{}\mathsf{fresh}\}$ gives the same; validate: passes. \code{foldl} in \code{core/List} is this shape with its callback's class carrying the accumulator's ownership ($\mathit{supply}\ \mathsf{owned},\allowbreak{}\mathit{need}\ \mathsf{owned}$), so \code{List.\allowbreak{}foldl xs [] (λx acc → List.\allowbreak{}push acc x)} is accepted and runs as research 42's R0 did (2–7× faster than the trie on builds, §6.4 there). \textbf{Solved.} Mutual recursion is the same fixpoint over the group. A group that reaches the round bound (it cannot: the lattice is finite) would be a $\top$ summary and a limit report.

\phantomsection\bheading{subsection}{5.6\quad Branches that keep the old value}{subsection}{\protect\numberline{5.6}Branches that keep the old value}

\begin{Verbatim}
case List.get xs i of
    Just v  → v
    Nothing → List.set xs i d ▷ List.get … 
\end{Verbatim}

$v$ is an element, not the cell; the \code{set} writes a slot. Per-arm liveness: \code{xs} is live after the \code{set} only on the \code{Nothing} arm where it is the operand. Accepted. The shape Rust's problem case \#3 could not accept for eight years (research 68 §7.1) does not arise because beni has no references into a list: \code{get} returns the element, not a pointer to the slot.

\begin{Verbatim}
ys = if c then List.set xs 0 1 else xs
\end{Verbatim}

Both arms: $A(\mathit{ys}) = \{\mathit{site}(\mathit{xs}\text{'s cell})\}$ either way — in the then-arm as the demand's result, in the else-arm as an alias. The demand in the then-arm asks whether \code{xs} is live after it on that arm: it is not (the else-arm's read of \code{xs} is not on the then path). Accepted; \code{ys} owns the cell in both arms. \textbf{Solved}: Futhark's union-of-both-branches rule (PLDI 2017 ALIAS-IF) is kept for \emph{cells} (the result may be either) but liveness is per path, so the false error of a value “used” in the arm not taken does not happen.

What is still coarse: \code{xs = if c then a else b; List.\allowbreak{}set xs 0 1} demands both $a$ and $b$ unique, so a later read of $a$ on the path where \code{xs = b} is a false error. A known imprecision; §9 counts it.

\phantomsection\bheading{subsection}{5.7\quad \code{foreign} and platform boundaries}{subsection}{\protect\numberline{5.7}\code{foreign} and platform boundaries}

A \code{foreign} has no body. Its declaration carries the summary as it carries the rung and \code{sync} (\code{boundary.md} §4). \textbf{The default is pessimistic}: a \code{foreign} parameter that says nothing is $\mathsf{shared}$ (the sibling may keep it — \code{boundary.md} §4 promises only that a sibling “never writes a list it was given, never keeps one it will later write”, which allows keeping one to \emph{read}, as \code{Hosted}'s key storage does) and a result that says nothing is $\top$. A sibling that only reads its argument during the call says $\mathsf{borrowed}$; one that returns “a fresh plain array that nothing else holds” (the existing rule) says $\mathsf{fresh}$; \code{Debug.log}'s result says $=\pi _{2}$. So a platform author who writes nothing loses precision and never soundness, and the words are what \code{boundary.md} §4's recipe gains. \code{Js.from x} marks $x$ escaped forever (research 42's pin, O7). A \code{Js.call} that passes a list is a \code{foreign} call with no summary: escaped. The checks of \code{boundary.md} §4 are unchanged; the new promise is tested the way the rung is: by review and by the \code{run/} fixtures that pin each row (A8). \textbf{Solved}, by contract; the residual risk is a sibling that lies, which is a correctness bug in privileged code, the same class as a wrong rung.

\phantomsection\bheading{subsection}{5.8\quad Fibers and commands capturing values}{subsection}{\protect\numberline{5.8}Fibers and commands capturing values}

\code{Cmd.\allowbreak{}perform (λsend → send (Saved model.\allowbreak{}todos))} captures \code{model.todos} in a thunk that the runtime stores and runs later (\code{boundary.md} §9.8.2). \code{Cmd.perform}'s summary: its argument $\mathit{escapes}=\mathsf{yes}$ — so every capture of the lambda is escaped at the capture, and \code{model.todos}'s cell is shared from then on, in this arm \emph{and in every later dispatch that receives the same cell in its model} (§5.9). \code{Task.spawn}, \code{Task.spawnIn}, \code{Task.par}'s thunks, \code{Ref.set}, \code{Queue.offer}: the same. A thunk the callee calls synchronously and does not keep (\code{Task.scope}'s body, \code{Task.bracket}'s acquire) is $\mathit{calls}\le 1,\allowbreak{}\mathit{escapes}\ \mathsf{no}$ and captures nothing permanently — except \code{bracket}'s \emph{release}, which is stored on the fiber's finaliser list and is $\mathit{escapes}=\mathsf{yes}$. These facts are declarations on \code{core/Task}'s \code{foreign}s and inferred for its beni wrappers. \textbf{Solved} by escape; the cost is that a list handed to any command or fiber is shared until the end of the program, because the analysis cannot see when the fiber reads it. The message names the capture (§6.2, example 3).

\phantomsection\bheading{subsection}{5.9\quad The TEA \code{update} on \code{browser-direct}: the handoff protocol}{subsection}{\protect\numberline{5.9}The TEA \code{update} on \code{browser-direct}: the handoff protocol}

Today's \code{browser} cannot host the rule: the runtime “keeps the rows it rendered (\code{s.b})” and “every patched hole keeps its last value” (research 42 §9), so every rendered list is a content holder the program cannot see, and a \code{List.\allowbreak{}push model.\allowbreak{}rows x} in \code{update} is an error on every message. Research 42 §6.1 measured it: with the runtime as written “R1 and R2 never write the table in place”. \code{browser-direct} removes the holds by design: no \code{view} at run time, no \code{lastRendered}, slots hold leaf values (§5.3), instance arrays hold \emph{items} and compare them by identity (§6.1–§6.2), scripts are guarded by tags and never by list identity (§6.2), and “the old model is dead after the arm” (§7.4). What remains is to state the contract the checker reads.

\textbf{The protocol (P1–P6), as obligations on the platform's lowering and runtime:}

\begin{itemize}
\item \textbf{P1 — handover.} The handler calls the arm with \code{model} \textbf{consumed}: after the arm, the runtime holds no content reference to any list reachable from the old model. The module-level \code{model} variable is reassigned by the arm (\code{model = …}, §4.1 step 1).
\item \textbf{P2 — reads before the arm are the handler's.} Anything the handler needs from the old model (§4.1 step 0: index symbols, a derived value's old input) is read \emph{before} the arm into locals; a list read this way is an ordinary holder of the handler, analysed by the same rules (so a script that reads the old list and compares it after the arm is a sharing error the lowering must not write — today's scripts read scalars and the \emph{new} list, §6.2).
\item \textbf{P3 — items are read only by the lowering's own code, and only as the lowering says.} \code{insts[i].it} is compared by \code{===} by the scripts (A4), and a row's groups read the \emph{new} item. The one program-visible read through an instance is a \textbf{listener body}: \code{browser-direct} §4.2 reads a handler's arguments “at the event, not captured” — \code{onClick=\{Select model.\allowbreak{}id\}} reads \code{model.id} when the click comes, and a row's handler reads “a row's instance through its root”. A listener body is therefore program code the checker analyses: it reads the model and the item \emph{before} the arm (it computes the payload), and what it hands the arm is P6's.
\item \textbf{P4 — slots are leaves or recomputed.} A derived value's slot holding a list is recomputed by every key whose write set conflicts with its reads (§5.4) before any group reads it; the old slot value is never dereferenced after the arm. This is \code{write-sets.md} §1.3's soundness, which already covers in-place writes once “changed” is read as “possibly changed in content” rather than “not \code{===}” (§7.4 there; §7.4 here).
\item \textbf{P5 — the entry state is a whole-program fixpoint of the checker's own.} \code{update} has no caller to discharge K2–K4 for it, so the checker computes the abstract value of \code{model} at entry itself: the join, over \code{init}'s result and every arm's result, of the cell sets at every model path, iterated to a fixpoint over the arms (one round per cell the arms can place, in practice two). Two things come out of it. \textbf{Aliasing between paths}: \code{init = xs = List.\allowbreak{}range 1 3; \{ a = xs, b = xs \}} puts one cell at \code{.a} and \code{.b}, so a \code{push} on \code{model.a} fails K4 against $(\mathit{model},\allowbreak{}.b)$ — the write-set pass cannot see this (its \code{Rec(none, …)} places two \code{Lst} nodes with no alias fact) and the checker must. \textbf{Escapes across dispatches}: a cell any arm captures in a \code{Cmd}, a \code{Sub}, a stored closure, a \code{Debug}/\code{Js} call, a \code{Send}, or a message payload a view event reads (P6) is escaped at every entry, because the runtime may hand the same cell back in the next model. \code{browser-direct} §7.4's condition (2) names the same set, but its scan is S7's future work (§11 item 4 there) and \code{write-sets.md} §3.3 analyses no command (“the command is not analysed”); the checker owns this computation. The entry summary is then $\mathit{use}(\mathit{model},\allowbreak{}p)=\mathsf{shared}$ for every escaped or aliased path, $\mathsf{consumed}$ otherwise.
\item \textbf{P6 — payloads.} A message's payload list is unique at the arm \textbf{only if nothing else holds it when the arm runs}. Two senders exist. A program sender — a fiber calling \code{send (Got rows)} — gives it up: \code{Send msg} is asserted $\mathsf{consumed}$ for the message, so a sender that reads \code{rows} afterwards is rejected at the sender, and the arm may edit \code{rows} in place. A \textbf{view event} whose handler expression is a model path or a row item (\code{onClick=\{Keep model.\allowbreak{}rows\}}, \code{onClick=\{Tag row.\allowbreak{}tags\}}) hands the arm a cell the model still holds: such a payload path is \textbf{shared at the arm} (it is one of P5's escapes), and the message offers to copy it in the arm or to send an id instead of the list.
\end{itemize}

Under P1–P6 the derivation for the table app's \code{Swap} and \code{Update} and for the browser corpus's \code{KeyedInPlace} (\code{List.\allowbreak{}update model.\allowbreak{}rows 1 …}, \code{List.\allowbreak{}push model.\allowbreak{}rows …}) is: \code{model} consumed at entry; \code{model.rows} not escaped by any arm (nothing captures it); the arm's operand \code{model.rows} has holders $(\mathit{model},\allowbreak{}.\mathit{rows})$ and nothing else; \code{model} is dead after the arm (P1); accepted, and the write is \code{a[i] = v} on the plain array — what \code{browser-direct} §7.4 calls “container-owned” and planned for S7, derived here from the general rule rather than from a model-specific one. \textbf{Solved}, conditional on P1–P6 being written into \code{browser-direct.\allowbreak{}md} and kept by its lowering; P2 and P5 are the obligations that need sentences the document does not have, and P5 is work: the entry fixpoint is the one whole-program computation in this design.

\textbf{Derived values and \code{verifyAll} are reads, never demands.} A \code{let} of \code{view} compiled into a derived value (\code{browser-direct} §5.4) — TodoMVC's \code{<For each=\{List.\allowbreak{}filter model.\allowbreak{}todos (visible model.\allowbreak{}filter \_\allowbreak{})\}>} is one — is recomputed by handlers while the page's \code{model} is live for the whole dispatch and beyond; and §8.3's development \code{verifyAll()} re-evaluates every derived value after each dispatch “and must change nothing unless the compiler is wrong”. So no expression a derived value or \code{verifyAll} evaluates may be a demand: the rebuilds there are the fresh-building forms. Under the rule as asked that is automatic (\code{filter} and \code{map} are not demands); under §10 change 1 it must be said, and it is the reason change 1 cannot be one function (§10).

\phantomsection\bheading{subsection}{5.10\quad \code{List.copy}}{subsection}{\protect\numberline{5.10}\code{List.copy}}

\code{copy : List a [borrowed] → List a [fresh]} — a shallow copy, \code{a.slice()} (or \code{\$plain().slice()} of a view). Shallow: the elements are the same values; records are immutable and need no copying; a \emph{nested} list inside an element is still shared, and editing it is still an error whose message names the inner holder, with the fix \code{List.\allowbreak{}map xs List.\allowbreak{}copy} (or a copy at the inner edit). The identity of a copy is new, so a \code{For} keyed by reference sees new rows — which is the right answer for a list whose contents the program is about to change behind the original's back. \code{copy} of a list already unique is legal and wasteful; the report mode counts such calls so that the owner can see whether a \code{copy\_\allowbreak{}of\_\allowbreak{}unique} warning is wanted (rule 7: a warning, never a refusal). \textbf{Solved.}

\phantomsection\bheading{subsection}{5.11\quad Views: a pattern's \code{rest}, \code{tail}, \code{drop}}{subsection}{\protect\numberline{5.11}Views: a pattern's \code{rest}, \code{tail}, \code{drop}}

\code{[ x, …rest ]} binds \code{rest} to a view $(b,\allowbreak{}o)$ of \code{xs}'s base (\code{backend.md} §4). A view is a holder of the base cell, and a write through a view is a write to the base: \code{List.\allowbreak{}set rest i v} is \code{b[o + i] = v}. So the demand on \code{rest} asks that every holder of $b$ — \code{xs} included — be dead. In the common walk (\code{go rest …} with \code{xs} not read again) it is; where \code{xs} is read after, it is a correct rejection: the elements \emph{are} shared. A view is always a suffix of its base (\code{backend.md} §4), so \code{push} on a view is \code{b.push(v)} and a new header, O(1); \code{pop} likewise. The scalar-view rewrite (\code{backend.md} §8) removes most views from loops anyway. \textbf{Solved.}

\phantomsection\bheading{subsection}{5.12\quad \code{Ref}, \code{Queue} and other mutable cells}{subsection}{\protect\numberline{5.12}\code{Ref}, \code{Queue} and other mutable cells}

\code{Ref.get r} returns an alias of whatever is in the cell; \code{Ref.set r xs} stores \code{xs}; the analysis cannot see who else called \code{Ref.get}. So a list stored in a \code{Ref} is \textbf{escaped}, every \code{Ref.get} result is $\top$-held, and \code{Ref.\allowbreak{}update r (λxs → List.\allowbreak{}push xs x)} is rejected with a limit message (the checker cannot see the cell's other readers). \textbf{Open as a precision limit}: a linear-cell discipline (the cell is the only holder and \code{update}'s callback consumes its argument) would need a proof that no \code{Ref.get} result is live across the \code{update}, which is a whole-program fact about an unbounded set of readers. The fix is \code{List.copy} in the callback, or a \code{Ref} of a \code{Dict}-like persistent structure. Measured by §9.

\phantomsection\bheading{subsection}{5.13\quad Messages as values, \code{Html.map}, subscriptions}{subsection}{\protect\numberline{5.13}Messages as values, \code{Html.map}, subscriptions}

A message stored in the model, logged whole or mapped by \code{Html.map} with a non-constructor function is a value (\code{browser-direct} §4.4); its payload lists are shared where the message is shared. A handler payload passed as parameters (the common case) is P6. \textbf{Solved.}

\phantomsection\bheading{subsection}{5.14\quad An undo history, a diff, a snapshot}{subsection}{\protect\numberline{5.14}An undo history, a diff, a snapshot}

\code{\{ model | history = [ model.\allowbreak{}todos, …model.\allowbreak{}history ], todos = List.\allowbreak{}push model.\allowbreak{}todos t \}}: the record update evaluates fields in written order (\code{language.md} §6), so \code{history} holds the old cell when the \code{push} runs — a live holder; \textbf{rejected}, correctly, and the message says that the undo entry would change with the push and offers \code{List.\allowbreak{}copy model.\allowbreak{}todos} in the history (O(n) per edit, the price of a version). A persistent-vector library in beni (the trie, as a \code{Dict}-like module with no in-place writes) would give O(log n) versions for programs that want them; it is a library, not a representation of \code{List}. \textbf{Solved}, with the cost stated.

\phantomsection\bheading{subsection}{5.15\quad Debug, the crash screen, \code{Debug.toString}}{subsection}{\protect\numberline{5.15}Debug, the crash screen, \code{Debug.toString}}

\code{Debug.log tag x} reads $x$ now and returns it (\code{write-sets.md} §3.7 row); \code{Debug.toString} reads now. Neither is a holder after the call. The development crash screen must print only the exception and the stack (A9); if it ever prints the model, it reads a half-updated record and the two semantics differ \emph{after} the program has stopped — a difference in a debugging aid, not in the program, but one to state. \textbf{Solved}, with A9 recorded.

\phantomsection\bheading{subsection}{5.16\quad The release optimiser's single-use inlining and specialisation}{subsection}{\protect\numberline{5.16}The release optimiser's single-use inlining and specialisation}

Inlining substitutes a body at a call, which preserves evaluation order and the holders (\code{language.md} §6); the verdict was computed on the unspecialised source and is not recomputed. Whole-program specialisation (\code{backend.md} §9) may make a copy of a callee per site, which may be \emph{more} unique than the general summary; the rule ignores it, so the accepted set is the source's, never the optimiser's (McCall). \textbf{Solved} by construction, with the cost that a program the specialiser could have proved safe is still rejected.
